% !TeX root = GeoDynamics.tex
\section{Schedule-Wise Contact Resolution}

This section documents the current schedule-wise contact combination rule. The
goal is to make the Python path match the GPU execution model as closely as
possible before porting the dynamics to GLSL.

The rule is source-owned:
\begin{itemize}
    \item one source particle computes its own velocity update,
    \item the source particle may read target particles and wall contact state,
    \item the source particle does not write target particle velocity,
    \item pair contacts, wall contacts, and final escape correction are applied
          in a fixed schedule.
\end{itemize}

This replaces the older cluster solve for the GLSL-portable path. The cluster
solver is retained only as a diagnostic comparator.

\subsection{Definitions}

For a source particle \(i\), let its current velocity be
\[
    \mathbf{v}_i .
\]

For each active particle contact \(c\), the contact resolver updates the
contact state and returns a source-owned contact momentum vector:
\[
    \Delta \mathbf{p}_{i,c}.
\]
This vector is computed from the contact's geometric/internal momentum scalar
and the contact normal.  The scheduler sums these momentum vectors and applies
one source velocity update.

\subsection{Algorithm}

For each source particle \(i\), the current schedule-wise algorithm is:

\begin{enumerate}
    \item Start with the particle's current velocity:
    \[
        \mathbf{v}_{base} = \mathbf{v}_i .
    \]

    \item Initialize the pair-contact momentum accumulator:
    \[
        \Delta \mathbf{p}_{sum} = (0,0).
    \]

    \item For each active particle-particle contact owned by source particle
    \(i\):
    \begin{enumerate}
        \item Update the source-owned contact state.

        \item If this is the first active pair contact, set:
        \[
            \mathbf{v}_{base} = \mathbf{v}_{i,c}^{0}
        \]
        where \(\mathbf{v}_{i,c}^{0}\) is the source velocity at first contact.

        \item Compute the contact momentum vector:
        \[
            \Delta \mathbf{p}_{i,c}
            =
            -p_{i,c}\mathbf{n}_{i,c}.
        \]

        \item Accumulate the contact momentum vector:
        \[
            \Delta \mathbf{p}_{sum}
            \leftarrow
            \Delta \mathbf{p}_{sum}
            +
            \Delta \mathbf{p}_{i,c}.
        \]
    \end{enumerate}

    \item If there are active pair contacts, apply the summed momentum once:
    \[
        \mathbf{v}_{pair}
        =
        \mathbf{v}_{base}
        +
        \frac{\Delta \mathbf{p}_{sum}}{m_i}.
    \]
    If there are no pair contacts, use:
    \[
        \mathbf{v}_{pair} = \mathbf{v}_i .
    \]

    \item Apply wall contacts sequentially to \(\mathbf{v}_{pair}\). For each
    active wall contact \(w\), compute the wall-predicted normal velocity and
    replace only the normal component of the current accumulator:
    \[
        \mathbf{v}_{wall}
        =
        \mathbf{v}_{acc}
        +
        \left(
            v_{n,w}^{\ast}
            -
            \mathbf{v}_{acc}\cdot\mathbf{n}_w
        \right)
        \mathbf{n}_w .
    \]
    After each wall contact:
    \[
        \mathbf{v}_{acc} \leftarrow \mathbf{v}_{wall}.
    \]

    \item Apply the pair rebound escape scheduler. For each active rebound pair
    contact, compute the desired outward relative normal velocity. If the
    current relative normal velocity is too small, adjust only the source
    particle accumulator:
    \[
        \mathbf{v}_{acc}
        \leftarrow
        \mathbf{v}_{acc}
        -
        \frac{1}{2}
        \left(
            u_{desired}
            -
            u_{current}
        \right)
        \mathbf{n}_{ij}.
    \]

    \item Write the final accumulator back to the source particle:
    \[
        \mathbf{v}_i \leftarrow \mathbf{v}_{acc}.
    \]
\end{enumerate}

\subsection{Pseudocode}

\begin{verbatim}
for each source particle i:
    base_velocity = current_velocity(i)
    accumulated_momentum = (0, 0)
    has_pair = false

    for each active particle contact c owned by i:
        contribution = pair_momentum_contribution(i, c)
        if contribution is invalid:
            continue

        if not has_pair:
            base_velocity = contribution.base_velocity
            has_pair = true

        accumulated_momentum += contribution.contact_momentum_vector

    if has_pair:
        pair_velocity =
            base_velocity + accumulated_momentum / mass(i)
    else:
        pair_velocity = current_velocity(i)

    final_velocity = pair_velocity

    for each active wall contact w owned by i:
        final_velocity = wall_prediction(i, w, final_velocity)

    for each active rebound pair contact c owned by i:
        final_velocity =
            apply_pair_escape_if_needed(i, c, final_velocity)

    velocity(i) = final_velocity
\end{verbatim}

\subsection{Known Question}

The current combination rule sums source-owned contact momentum vectors.  This
matches the geometric multi-contact convention more closely than averaging
velocity predictions, but it is still not a coupled global solve.  Dense contact
groups may still require limiting, weighting, or another source-owned
geometric rule.
